% !PRINTED 458
\weberpage{458}

Besteht umgekehrt eine solche Gleichung, so folgt durch Vergleichung
mit (4) in der früher angewandten Art, wenn $x, y$ ganze Zahlen sind:

\weq{7}{
\begin{aligned}
\partial\beta &= Ax, \\
c\alpha-a\gamma &= Cx, \\
c\beta+\partial\alpha-a\delta &= Bx, \\
c\beta-\partial\alpha-a\delta &= y,
\end{aligned}
}

woraus:

\weq{8}{
\begin{aligned}
2\partial\alpha &= Bx-y, \\
2(c\beta-a\delta) &= Bx+y, \\
4n &= y^2-Dx^2.
\end{aligned}
}

Setzen wir $A$ relativ prim zu $3n$ voraus, so muß $\partial = 1$,
$a = n$ angenommen werden; denn nach (7) muß unter dieser
Voraussetzung $x$ durch $\partial$ teilbar sein, und folglich auch
$c\alpha-a\gamma$ und $c\beta-a\delta$, woraus folgt, daß auch $a, c$
durch $\partial$ teilbar sein müßten, während doch $a, c, \partial$
keinen gemeinsamen Teiler haben. Also ist

\weq{9}{
\begin{aligned}
2\alpha &= Bx-y, & \beta &= Ax, \\
2n\gamma &= c(Bx-y)-2Cx, & 2n\delta &= 2cAx-Bx-y.
\end{aligned}
}

Die Zahlen $x, y$ sind hier außer den in \S\,114, 1), 2) angegebenen
Beschränkungen wegen (3) noch der unterworfen, daß
$y^2-Dx^2 \equiv -1\ (\mathrm{mod}\ 3)$ werde. \gesperrt{Dadurch ist
ausgeschlossen, daß $D$ durch 3 teilbar sei.} $x$ muß von Null
verschieden sein, und wir können es, ohne die Allgemeinheit zu
beschränken, positiv annehmen. (Eine gleichzeitige
Vorzeichenänderung von $x$ und $y$ bewirkt nur eine
Vorzeichenänderung von $\alpha, \beta, \gamma, \delta$, die ohne
Belang ist.)

Ist $D \equiv -1\ (\mathrm{mod}\ 3)$, so müssen $x$ und $y$ durch 3
unteilbar sein; ist $D \equiv 1\ (\mathrm{mod}\ 3)$, so ist $y$ durch
3 teilbar, $x$ nicht teilbar. Im übrigen können die Zahlen $x, y$
beliebig angenommen werden. Es ist dann $c$ bestimmt aus den beiden
(miteinander verträglichen) Kongruenzen:

\[
\begin{aligned}
cAx &\equiv \frac{Bx+y}{2} \\[4pt]
c\,\frac{Bx-y}{2} &\equiv Cx
\end{aligned}
\quad (\mathrm{mod}\ n)
\]

und kann, da $n$ durch 3 unteilbar ist, noch der Bedingung

\[
c \equiv 0\ (\mathrm{mod}\ 3)
\]

unterworfen werden.
